\documentclass[10pt,landscape]{article}

\usepackage[margin=0.32in]{geometry}
\usepackage[T1]{fontenc}
\usepackage[utf8]{inputenc}
\usepackage{lmodern}
\usepackage{textcomp}
\usepackage{microtype}
\usepackage[table,dvipsnames]{xcolor}
\usepackage{array}
\usepackage{booktabs}
\usepackage{enumitem}
\usepackage{hyperref}
\usepackage{listings}
\usepackage{multicol}
\usepackage{tabularx}
\usepackage{titlesec}

\pagestyle{empty}
\setlength{\parindent}{0pt}
\setlength{\parskip}{1.5pt}
\setlength{\columnsep}{13pt}
\setlength{\columnseprule}{0.2pt}
\setcounter{secnumdepth}{0}

\definecolor{ToshBlue}{HTML}{0F4C81}
\definecolor{ToshInk}{HTML}{1F2937}
\definecolor{ToshMuted}{HTML}{5B6470}
\definecolor{ToshPanel}{HTML}{F4F7FB}
\definecolor{ToshRule}{HTML}{CBD5E1}
\definecolor{ToshCodeBg}{HTML}{F8FAFC}

\color{ToshInk}

\titleformat{\section}
  {\large\bfseries\color{ToshBlue}}
  {}
  {0pt}
  {}
\titlespacing*{\section}{0pt}{0.4em}{0.12em}

\setlist[itemize]{leftmargin=1.05em,itemsep=0.5pt,topsep=1.5pt,parsep=0pt,partopsep=0pt}

\lstdefinestyle{tosh}{
  basicstyle=\ttfamily\scriptsize,
  backgroundcolor=\color{ToshCodeBg},
  frame=single,
  rulecolor=\color{ToshRule},
  framerule=0.4pt,
  xleftmargin=0pt,
  xrightmargin=0pt,
  aboveskip=0.18em,
  belowskip=0.24em,
  columns=fullflexible,
  keepspaces=true,
  showstringspaces=false,
  breaklines=true,
  literate={°}{{\textdegree}}1
           {μ}{{\ensuremath{\mu}}}1
           {Ω}{{\ensuremath{\Omega}}}1
}
\lstset{style=tosh}

\newcommand{\sheettitle}[3]{
  {\LARGE\bfseries\color{ToshBlue} ToSh Cheat Sheet: #1\par}
  \vspace{0.08em}
  {\small\color{ToshMuted} #2 \hfill #3\par}
  \vspace{0.18em}
  {\color{ToshRule}\rule{\linewidth}{0.8pt}}
  \vspace{0.12em}
}

\newcommand{\note}[1]{
  \vspace{0.12em}
  \fcolorbox{ToshRule}{ToshPanel}{%
    \parbox{\dimexpr\linewidth-2\fboxsep-2\fboxrule\relax}{\footnotesize #1}%
  }
  \vspace{0.12em}
}

\newcommand{\tighttable}[2]{
  {\footnotesize
  \renewcommand{\arraystretch}{1.08}
  \begin{tabularx}{\linewidth}{@{}#1@{}}
  #2
  \end{tabularx}}
}


\begin{document}
\footnotesize
\sheettitle{Units, Dates, And Time}{StorageSize, TimeSpan, TemporalAmount, date math, and physical quantities}{spec + runtime behavior}

\begin{multicols*}{3}

\section{Temporal Literals}
\begin{lstlisting}
5s
100ms
2m30s
1w2d4h5m23s
1y2mo
2026-03-25
\end{lstlisting}

\tighttable{>{\ttfamily\raggedright\arraybackslash}p{0.30\linewidth}X}{
\toprule
Literal & Runtime type \\
\midrule
\texttt{2d}, \texttt{5s} & \texttt{TimeSpan} \\
\texttt{1y2mo} & \texttt{TemporalAmount} \\
\texttt{2026-03-25} & \texttt{DateTimeOffset} \\
\bottomrule
}

\begin{lstlisting}
date now
date parse 2026-03-23T14:05:00Z
date date-only 2026-03-23T14:05:00Z
date time-only 2026-03-23T14:05:00Z
timespan 1w2d
\end{lstlisting}

\section{Date Arithmetic}
\begin{lstlisting}
var now = (date now)
var yesterday = (($now) - 1d)
var nextWeek = (($now) + 1w)
\end{lstlisting}

\begin{lstlisting}
var jan31 = (date 2026-01-31)
var febEnd = (($jan31) + 1mo)
echo $febEnd.Month $febEnd.Day
\end{lstlisting}

\note{Calendar units like \texttt{mo} and \texttt{y} produce \texttt{TemporalAmount}; pure clock units produce \texttt{TimeSpan}.}

\section{Storage Sizes}
\begin{lstlisting}
1kb
10mb
2.5gb
\end{lstlisting}

\begin{lstlisting}
ls -la | sort Size
ls -la | where _.Size >= 1kb | get Name
ls -la | where _.Size? > 1000 | get Name
\end{lstlisting}

\begin{itemize}
\item \texttt{StorageSize} values compare naturally in pipelines.
\item They can be compared against either numbers or size literals.
\item Display is configurable from \texttt{\$tosh.Config.Display}.
\end{itemize}

\columnbreak

\section{Quantity Literals}
\begin{lstlisting}
100`m
5`kg
9.8`m/s^2
2.5`kWh
72`degF
90`deg
72°F
90°
\end{lstlisting}

\tighttable{>{\ttfamily\raggedright\arraybackslash}p{0.28\linewidth}X}{
\toprule
Piece & Meaning \\
\midrule
\texttt{100`m} & number + backtick + unit expression \\
\texttt{m/s\^{}2} & division and exponent are allowed \\
\texttt{km}, \texttt{MB} & SI prefixes resolve automatically \\
\texttt{90°}, \texttt{72°F} & degree-sign shorthand for angle / temperature \\
\bottomrule
}

\begin{lstlisting}
100`m + 50`m
100`km + 500`m
100`m / 20`s
5`kg * 9.8`m/s^2
\end{lstlisting}

\note{Quantities enforce dimensional analysis: compatible linear quantities add/compare cleanly; incompatible dimensions raise an error. Absolute temperatures convert and compare, but arithmetic awaits explicit temperature-difference units.}

\section{Convert Units}
\begin{lstlisting}
2`mi as `ft
2`hr as `s
10`m/s as `mph
20°C as `°F

$distance.To("km")
\end{lstlisting}

\note{The backtick after \texttt{as} marks a unit target. Without it, \texttt{as Length} remains an ordinary type conversion.}

\section{Named Quantity Categories}
\tighttable{>{\ttfamily\raggedright\arraybackslash}p{0.30\linewidth}X}{
\toprule
Category & Example units \\
\midrule
Length & \texttt{m}, \texttt{km}, \texttt{ft}, \texttt{mi} \\
Mass & \texttt{kg}, \texttt{g}, \texttt{lb} \\
Duration & \texttt{s}, \texttt{ms}, \texttt{min}, \texttt{hr} \\
Temperature & \texttt{K}, \texttt{degC}, \texttt{degF} \\
DataSize & \texttt{B}, \texttt{kB}, \texttt{MB}, \texttt{GB} \\
Speed & \texttt{m/s}, \texttt{km/h}, \texttt{mph} \\
Energy & \texttt{J}, \texttt{kJ}, \texttt{kWh}, \texttt{BTU} \\
Angle & \texttt{rad}, \texttt{deg} \\
\bottomrule
}

\section{Quantity Members}
\begin{lstlisting}
var d = 5`km
echo $d.Magnitude
echo $d.UnitSymbol
echo $d.BaseValue
echo $d.CategoryName
\end{lstlisting}

\begin{itemize}
\item \texttt{Magnitude}: numeric value in the original unit
\item \texttt{UnitSymbol}: unit text as entered
\item \texttt{BaseValue}: SI-normalized value
\item \texttt{CategoryName}: named physical category when recognized
\end{itemize}

\columnbreak

\section{Temporal Pipelines}
\begin{lstlisting}
func recentFiles(age: TimeSpan) {
    ls -la | where _.Type == file | where _.Modified > ((date now) - $age)
}

recentFiles 7d | get { Name, Modified, Size }
\end{lstlisting}

\begin{lstlisting}
date parse 2026-03-23T14:05:00Z -d -t
date parse 2026-03-23T14:05:00Z | cast dateonly
date parse 2026-03-23T14:05:00Z | cast timeonly
\end{lstlisting}

\section{Date Helpers}
\begin{lstlisting}
date today
date tomorrow
date utc-now
date from-unix 1742738700
\end{lstlisting}

\section{Other Intrinsic Literals}
\begin{lstlisting}
127.0.0.1
::1
\end{lstlisting}

\begin{itemize}
\item Bare IPv4 and IPv6 values parse as \texttt{System.Net.IPAddress}.
\item Typed parameters can coerce string input into \texttt{ipaddress} values.
\end{itemize}

\section{Practical Examples}
\begin{lstlisting}
# files newer than two days ago
ls -la | where _.Modified > ((date now) - 2d)

# average speed for a run
echo (5`km / 24`min)

# filter by size
find . -name *.log | where _.Size >= 10mb
\end{lstlisting}

\section{Display Tuning}
\begin{lstlisting}
view datetime table unix
view duration table seconds
view size bytes
view dateonly scalar relative
\end{lstlisting}

\note{\textbf{Bridge behavior:} ToSh can promote plain \texttt{TimeSpan} and \texttt{StorageSize} values when they interact with the unit system, so date math and size comparisons stay ergonomic in real shell pipelines.}

\end{multicols*}
\newpage

\sheettitle{Units, Dates, And Time}{Engineering calculations, bridge types, and calendar-aware shell recipes}{page 2 / front-back reference}

\begin{multicols*}{3}

\section{Bridge Types}
\tighttable{>{\ttfamily\raggedright\arraybackslash}p{0.31\linewidth}X}{
\toprule
Form & Meaning \\
\midrule
\texttt{5s}, \texttt{2h30m} & \texttt{TimeSpan} clock durations \\
\texttt{1y2mo} & \texttt{TemporalAmount} calendar duration \\
\texttt{10mb} & \texttt{StorageSize} byte-oriented size \\
\texttt{5`km} & \texttt{Quantity} with dimensional analysis \\
\texttt{127.0.0.1} & \texttt{System.Net.IPAddress} literal \\
\bottomrule
}

\begin{itemize}
\item Use storage literals for files and memory, unit literals for physics and engineering math.
\item \texttt{date add/sub} accepts calendar-aware \texttt{TemporalAmount} values and fixed \texttt{DurationQuantity} values; \texttt{1mo} and \texttt{2`hr} keep their distinct meanings.
\item \texttt{sum} and \texttt{avg} handle numeric, \texttt{Quantity}, \texttt{StorageSize}, and \texttt{TimeSpan} streams directly.
\end{itemize}

\section{Derived Quantity Families}
\tighttable{>{\ttfamily\raggedright\arraybackslash}p{0.33\linewidth}X}{
\toprule
Category & Examples \\
\midrule
Area / Volume & \texttt{m\^{}2}, \texttt{acre}, \texttt{L}, \texttt{gal} \\
Acceleration / Force & \texttt{m/s\^{}2}, \texttt{N}, \texttt{lbf} \\
Power / Pressure & \texttt{W}, \texttt{kW}, \texttt{Pa}, \texttt{psi} \\
Charge / Resistance & \texttt{C}, \texttt{Ah}, \texttt{V}, \texttt{ohm} \\
Density / FlowRate & \texttt{kg/m\^{}3}, \texttt{L/min}, \texttt{gal/min} \\
Angle / AngularVelocity & \texttt{deg}, \texttt{rad}, \texttt{deg/s} \\
\bottomrule
}

\section{Dimension Math}
\begin{lstlisting}
var force = (12`kg * 9.81`m/s^2)
var work = ($force * 3`m)
var power = ($work / 10`s)

echo $force
echo $work
echo $power
\end{lstlisting}

\begin{lstlisting}
echo (60`mph > 25`m/s)
echo (7850`kg/m^3 * 0.01`m^3)
echo (2`L / 30`s)
\end{lstlisting}

\columnbreak

\section{Calendar-Aware Recipes}
\begin{lstlisting}
date today
date tomorrow
date yesterday
date utc-now
date from-unix-ms 1742738700000
\end{lstlisting}

\begin{lstlisting}
date today | date add 1mo
date 2026-03-25 | date sub 2w
date parse 2026-03-23T14:05:00Z | date time-only
\end{lstlisting}

\begin{itemize}
\item Month/year amounts clamp naturally at month ends.
\item Pipe a date into \texttt{date add} / \texttt{date sub} when you want calendar logic without extra variables.
\item Use \texttt{-d} and \texttt{-t} when you want both projections emitted at once.
\end{itemize}

\section{Temporal Filtering}
\begin{lstlisting}
ls -la |
where _.Type == file |
where _.Modified > ((date now) - 12h) |
sort Modified | reverse
\end{lstlisting}

\begin{lstlisting}
find . -name *.log |
where _.Modified < ((date now) - 30d) |
get { Name, Modified, Size }
\end{lstlisting}

\section{Size And Duration Aggregates}
\begin{lstlisting}
echo 1kb 2kb 3kb | sum
echo 1s 2s 3s | avg
ls -la | where _.Type == file | summarize --sum Size
\end{lstlisting}

\note{\texttt{sum} and \texttt{avg} accept compatible quantities and preserve the first value's display unit. \texttt{min} and \texttt{max} return the winning original value, including its unit. Summing absolute temperatures is intentionally rejected.}

\section{Display Profiles}
\begin{lstlisting}
view dateonly long
view timeonly 24h
view duration long
view size auto
\end{lstlisting}

\columnbreak

\section{Scientific And Ops Examples}
\begin{lstlisting}
# network throughput from bytes / time
echo (120`MB / 30`s)

# battery energy over time -> power
echo (2.5`kWh / 30`min)

# rotation rate
echo (180`deg / 2`s)
\end{lstlisting}

\begin{lstlisting}
# compare mixed units
echo (5`mi > 8`km)

# temperature math stays typed
echo 72`degF
echo 22`degC
\end{lstlisting}

\section{Working With Dates As Types}
\begin{lstlisting}
var dt = (date now)
var d = ($dt | cast dateonly)
var t = ($dt | cast timeonly)

echo $d
echo $t
\end{lstlisting}

\begin{lstlisting}
type-of (date now)
type-of 2d
type-of 1y2mo
type-of 5`km
\end{lstlisting}

\section{Practical File Retention}
\begin{lstlisting}
func staleFiles(age: TimeSpan, min: StorageSize) {
    ls -la |
    where _.Type == file |
    where _.Modified < ((date now) - $age) |
    where _.Size >= $min
}
\end{lstlisting}

\begin{lstlisting}
staleFiles 14d 100mb | get { Name, Size, Modified }
\end{lstlisting}

\section{Mental Model}
\begin{itemize}
\item \texttt{2d} is shell time, \texttt{2`m} is physical length, and \texttt{2mb} is storage.
\item Calendar durations (\texttt{mo}, \texttt{y}) are not just scaled seconds; they preserve calendar semantics.
\item Quantities keep dimensions honest, which makes shell calculations safer than raw doubles.
\end{itemize}

\end{multicols*}
\end{document}
